% ============================================================
%  PLANTILLA SYNTHECA, adaptada de la plantilla minimalista
%  del usuario (article) a clase BOOK.
%
%  Decisión de diseño: cada TEMA generado por /generar-sintesis es
%  un \chapter. Con el tiempo, una misma materia acumula capítulos
%  y el documento se convierte literalmente en tu propio libro de
%  esa materia; no hace falta un documento nuevo por tema.
%
%  Mismo criterio estético que la plantilla original: gris oscuro
%  (no negro puro), títulos centrados en mayúsculas con línea
%  divisoria, márgenes estrechos, sin decoración innecesaria.
%  Los entornos nuevos (bloque, ejemplo, chequeo, ejercicio) usan
%  el mismo principio de minimalismo visual que rige todo el
%  contenido: líneas finas para separar, nunca colores de fondo.
%  La única excepción es el entorno ejercicio, que lleva un
%  recuadro fino gris oscuro sin relleno: los ejercicios son la
%  parte del documento que el lector trabaja aparte (con lápiz y
%  papel), y el recuadro los separa visualmente de la teoría.
%
%  Regla tipográfica para el texto que va DENTRO de estos entornos
%  (la aplica redactor-pedagogico, ver plantilla-sintesis/SKILL.md):
%  puntuación normal del castellano. Sin rayas (---) como separador
%  de incisos, sin flechas unicode, sin listas de sustantivos con
%  dos puntos. Paréntesis, coma, punto y coma o punto seguido.
% ============================================================
\usepackage[utf8]{inputenc}
\usepackage[T1]{fontenc}
% NOTA: se remueve babel[spanish] por el mismo motivo que la
% plantilla original: el paquete de hifenación en español no está
% instalado en este entorno de compilación. Se renombran a mano los
% términos que babel normalmente traduciría (ver más abajo).
\usepackage{geometry}
\usepackage{amsmath}
\usepackage{amssymb}
\usepackage{amsthm}
\usepackage{enumitem}
\usepackage{tabularx}
\usepackage{booktabs}
\renewcommand{\familydefault}{\sfdefault}
\usepackage{titlesec}
\usepackage{fancyhdr}
\usepackage{ragged2e}
\usepackage[table]{xcolor}
\usepackage{array}
\usepackage{mdframed}

\geometry{
    top=2cm,
    bottom=2.2cm,
    left=2cm,
    right=2cm
}

% ---------- Colores (idénticos a la plantilla original) ----------
\definecolor{textgray}{RGB}{60, 60, 60}
\definecolor{rulegray}{RGB}{150, 150, 150}
\color{textgray}

% ---------- Metadatos editables por síntesis ----------
\newcommand{\miNombre}{Nombre Apellido}      % quién estudia (footer)
\newcommand{\materiaActual}{Materia}          % se resetea por materia
\newcommand{\setMateria}[1]{\renewcommand{\materiaActual}{#1}}

% ---------- Header / Footer (igual que el original) ----------
\pagestyle{fancy}
\fancyhf{}
\renewcommand{\headrulewidth}{0pt}
\renewcommand{\footrulewidth}{0pt}
\fancyfoot[C]{\small\color{textgray}\thepage}
\fancyfoot[R]{\footnotesize\color{textgray}\miNombre}
\fancyfoot[L]{\footnotesize\color{textgray}\materiaActual}

% ---------- Traducciones manuales (sin babel) ----------
\renewcommand{\tablename}{Tabla}
\renewcommand{\figurename}{Figura}
\renewcommand{\contentsname}{Índice}
\newcommand{\ok}{\checkmark}

% ============================================================
%  CAPÍTULO = TEMA. Mismo criterio visual que \section en la
%  plantilla original (centrado, mayúsculas, línea fina), pero un
%  punto más grande porque es el nivel jerárquico superior acá.
% ============================================================
\titleformat{\chapter}
  {\centering\color{textgray}\fontsize{18}{21}\selectfont\bfseries}
  {}
  {0pt}
  {}
  [\vspace{3pt}\color{rulegray}\titlerule\vspace{6pt}]
\titlespacing*{\chapter}{0pt}{0pt}{18pt}
% NOTA: a diferencia de \section (donde redefinir el comando funciona sin
% problema), redefinir \chapter directamente rompe el mecanismo interno de
% \tableofcontents en la clase book (corrompe la entrada de TOC). En vez de
% eso, se usa un comando auxiliar explícito: \temacapitulo{<Tema>} en lugar
% de \chapter{<Tema>} directo.
\newcommand{\temacapitulo}[1]{\chapter{\MakeUppercase{#1}}}

% ============================================================
%  SECTION = cada paso IMPRESO de la plantilla pedagógica
%  (Motivación, Marco teórico, Ejemplo resuelto, Errores comunes,
%  Ejercicios, y si hay herramienta, Traducción y Ejercicio guiado).
%  El paso 2 de la plantilla (conciliación de fuentes) es interno
%  al pipeline y no se imprime como sección.
%  Formato idéntico al de la plantilla original.
% ============================================================
\titleformat{\section}
  {\centering\color{textgray}\fontsize{13}{15}\selectfont\bfseries}
  {}
  {0pt}
  {}
  [\vspace{2pt}\color{rulegray}\titlerule\vspace{4pt}]
\titlespacing*{\section}{0pt}{18pt}{6pt}
\let\oldsection\section
\renewcommand{\section}[1]{\oldsection{\MakeUppercase{#1}}}

% ============================================================
%  SUBSECTION = micro-bloques dentro de "3. Marco teórico"
%  (3.a definición, 3.b prompt generativo; 3.c usa el entorno
%  \chequeo, no subsection, porque necesita numeración propia).
% ============================================================
\titleformat{\subsection}
  {\color{textgray}\normalsize\bfseries}
  {}
  {0pt}
  {}
\titlespacing*{\subsection}{0pt}{12pt}{4pt}

% ============================================================
%  ENTORNOS CUSTOM, uno por cada pieza funcional de la
%  plantilla pedagógica. Numerados por capítulo (=por tema).
% ============================================================

\theoremstyle{definition}
\newtheorem{bloqueteorico}{Bloque}[chapter]
% Paso 3.a: definición/fórmula del bloque.
% Uso normal:  \begin{bloque}{Distancia de Hamming} ... \end{bloque}
% El argumento opcional (etiqueta chica en gris) se reserva para el caso
% en que dos fuentes usan notación distinta o se contradicen y hace falta
% dejar claro cuál se sigue en ese bloque. No se usa por defecto.
% Uso excepcional: \begin{bloque}[notación de Lin y Costello]{Distancia de Hamming}
\newenvironment{bloque}[2][]{%
  \bloqueteorico%
  \noindent\textbf{\color{textgray}#2}%
  \ifx&#1&\else\ {\footnotesize\color{rulegray}(#1)}\fi%
  \par\vspace{2pt}\color{rulegray}\hrule\color{textgray}\vspace{4pt}%
}{%
  \vspace{4pt}
}

\theoremstyle{plain}
\newtheorem{ejemploresueltoct}{Ejemplo resuelto}[chapter]
% Paso 4: ejemplo resuelto. Argumento opcional: "completo" | "parcial"
% (alternar entre ambos formatos por regla de pedagogia-cognitiva).
% Si el formato se calibró contra una resolución real de la cátedra,
% cerrar el entorno con \citaEjercicio{...} (ver más abajo).
\newenvironment{ejemplo}[1][completo]{%
  \ejemploresueltoct%
  \ifthenelse{\equal{#1}{parcial}}%
    {\ {\footnotesize\itshape\color{rulegray}(parcial: completá el paso final)}}%
    {}%
  \par\vspace{2pt}\color{rulegray}\hrule\color{textgray}\vspace{4pt}%
}{%
  \vspace{4pt}
}

% Paso 3.c: chequeo de comprensión (adjunct questions). Va INMEDIATAMENTE
% después de cada bloque teórico, nunca concentrado al final del capítulo.
\newmdenv[
  linecolor=rulegray,
  linewidth=0.4pt,
  topline=true,
  bottomline=true,
  leftline=false,
  rightline=false,
  innertopmargin=6pt,
  innerbottommargin=6pt,
  skipabove=8pt,
  skipbelow=8pt
]{chequeobox}
\newenvironment{chequeo}{%
  \begin{chequeobox}%
  {\footnotesize\bfseries\color{rulegray} CHEQUEO RÁPIDO}\par\vspace{2pt}%
  \begin{enumerate}[leftmargin=1.4em, itemsep=2pt, label=\arabic*.]
}{%
  \end{enumerate}
  \end{chequeobox}
}

% ============================================================
%  Paso 6.a / 6.b (y 8): EJERCICIOS.
%  Recuadro fino gris oscuro, sin fondo, que separa cada ejercicio
%  del texto que lo rodea. Numerado por capítulo.
%
%  Uso:
%    \begin{ejercicio}[familia]{fase}
%      Enunciado...
%      \citaEjercicio{Similar al ejercicio 3 del Parcial 1 (2023)}   % opcional
%    \end{ejercicio}
%
%  - familia (opcional): se imprime en gris chico al lado del número,
%    para que quede trazable contra mapa-estudio.json.
%  - fase (obligatorio): bloque | intercalado | guiado. NO se imprime;
%    queda en el fuente para que rubric_check.py verifique el orden
%    (bloque antes que intercalado). La distinción visible la dan las
%    \subsection de la sección Ejercicios.
%  - \citaEjercicio: SOLO cuando el ejercicio se modeló sobre un
%    ejercicio real de skills/banco-ejercicios/ (guía, TP, parcial o
%    final). Se escribe en castellano natural, con la referencia
%    concreta: "Similar al ejercicio 4 de la Guía 3", "Adaptado del
%    ejercicio 2 del Parcial 2 (2024)". Nunca se inventa una cita: si
%    no hubo --estilo o el ejercicio no sale de uno real, no se pone.
% ============================================================
\definecolor{framegray}{RGB}{90, 90, 90}
\newmdenv[
  linecolor=framegray,
  linewidth=0.5pt,
  roundcorner=0pt,
  topline=true,
  bottomline=true,
  leftline=true,
  rightline=true,
  innertopmargin=7pt,
  innerbottommargin=8pt,
  innerleftmargin=10pt,
  innerrightmargin=10pt,
  skipabove=10pt,
  skipbelow=10pt
]{ejerciciobox}
\newcounter{ejercicionum}[chapter]
\renewcommand{\theejercicionum}{\thechapter.\arabic{ejercicionum}}
\newenvironment{ejercicio}[2][]{%
  \begin{ejerciciobox}%
  \refstepcounter{ejercicionum}%
  \noindent{\bfseries\color{textgray}Ejercicio \theejercicionum}%
  \ifx&#1&\else\ {\footnotesize\color{rulegray}[#1]}\fi%
  \par\vspace{4pt}\noindent\ignorespaces%
}{%
  \end{ejerciciobox}%
}
% Referencia al ejercicio real en el que se basa (opcional). Se coloca al
% final del enunciado, dentro de ejercicio o de ejemplo. Línea chica en
% itálica gris, alineada a la derecha.
\newcommand{\citaEjercicio}[1]{%
  \par\vspace{4pt}\noindent\hfill{\footnotesize\itshape\color{rulegray}#1}\par%
}

% Paso 5: errores comunes. Lista simple, sin caja (el tono neutro va
% en el TEXTO, no en la decoración, coherente con el minimalismo visual).
\newenvironment{erroresComunes}{%
  \begin{itemize}[leftmargin=1.4em, itemsep=4pt, label={\color{rulegray}\textbullet}]
}{%
  \end{itemize}
}

% Paso 1: motivación. Estilo narrativo en primera persona, apertura
% con oración temática (itálica, con leve sangría) antes del cuerpo.
\newenvironment{motivacion}{%
  \begin{quote}
  \itshape\color{textgray}
}{%
  \end{quote}
}

% Necesario para el \ifthenelse de \ejemplo
\usepackage{ifthen}

